% 集中定义模型图通用样式和通用组件
\tikzset{
  % 前向传播实线箭头
  dann edge/.style={
    -{Latex[length=1.5mm]}, % 箭头样式
    draw=teal!50!black, % 连线颜色
    line width=.55pt % 连线宽度
  },
  % 域适配/混淆损失虚线箭头
  dann confuse edge/.style={
    -{Latex[length=1.5mm]}, % 箭头样式
    draw=purple!65!black, % 连线颜色
    line width=.55pt, % 连线宽度
    % dashed % 虚线表示域混淆或分布对齐
  },
  % 输入数据节点基础样式
  dann data/.style={
    rounded corners=2pt, % 轻微圆角
    draw=blue!55!black, % 默认边框颜色
    fill=blue!9, % 默认填充颜色
    minimum width=18mm, % 最小宽度
    minimum height=10mm, % 最小高度
    align=center, % 居中换行
    font=\scriptsize % 使用较小字号
  },
  % Source数据节点样式
  dann source data/.style={
    dann data, % 继承数据节点基础样式
    draw=blue!75!black, % 源域边框色
    fill=blue!10 % 源域填充色
  },
  % Target数据节点样式
  dann target data/.style={
    dann data, % 继承数据节点基础样式
    draw=green!50!black, % 目标域边框色
    fill=green!12 % 目标域填充色
  },
  % 分类头节点样式
  dann head/.style={
    rounded corners=2pt, % 轻微圆角
    draw=orange!70!black, % 边框颜色
    fill=orange!13, % 填充颜色
    minimum width=18mm, % 最小宽度
    minimum height=9mm, % 最小高度
    align=center, % 居中换行
    font=\scriptsize % 使用较小字号
  },
  % GRL节点样式
  dann grl/.style={
    rounded corners=2pt, % 轻微圆角
    draw=purple!70!black, % 边框颜色
    fill=purple!12, % 填充颜色
    minimum width=16mm, % 最小宽度
    minimum height=9mm, % 最小高度
    align=center, % 居中换行
    font=\scriptsize % 使用较小加粗字号
  },
  % 域判别器节点样式
  dann domain/.style={
    rounded corners=2pt, % 轻微圆角
    draw=red!65!black, % 边框颜色
    fill=red!10, % 填充颜色
    minimum width=19mm, % 最小宽度
    minimum height=9mm, % 最小高度
    align=center, % 居中换行
    font=\scriptsize % 使用较小字号
  },
  % 输出节点样式
  dann output/.style={
    rounded corners=2pt, % 轻微圆角
    draw=gray!65, % 边框颜色
    fill=gray!8, % 填充颜色
    minimum width=17mm, % 最小宽度
    minimum height=8mm, % 最小高度
    align=center, % 居中换行
    font=\scriptsize % 使用较小字号
  },
  % 简化版共享CFE矩形样式
  dann cfe/.style={
    rounded corners=3pt, % 圆角
    draw=teal!60!black, % 边框颜色
    fill=teal!10, % 填充颜色
    minimum width=24mm, % 最小宽度
    minimum height=16mm, % 最小高度
    align=center, % 居中换行
    font=\scriptsize % 使用较小加粗字号
  },
  % Source EEG 输入组件
  pics/source eeg/.style={
    code={ % 定义Source EEG模块
      \eegDataCard{sourceData}{Source data}{dann source data} % 源域数据卡片
      \coordinate (source-east) at (sourceData.east); % 右侧连接点
    }
  },
  % Target EEG 输入组件
  pics/target eeg/.style={
    code={ % 定义Target EEG模块
      \eegDataCard{targetData}{Target data}{dann target data} % 目标域数据卡片
      \coordinate (target-east) at (targetData.east); % 右侧连接点
    }
  },
  % 简化版共享CFE组件
  pics/shared cfe/.style={
    code={ % 定义共享特征提取器
      \node[dann cfe] (sharedCFE) at (0, 0) {Shared CFE\\310$\to$64}; % 共享CFE矩形
      \node[font=\scriptsize, text=black] at (0, -.92) {Feature Extractor}; % CFE说明
      \coordinate (cfe-source-west) at ([yshift=2.5mm]sharedCFE.west); % 源域输入连接点
      \coordinate (cfe-target-west) at ([yshift=-2.5mm]sharedCFE.west); % 目标域输入连接点
      \coordinate (cfe-emo-east) at ([yshift=2.5mm]sharedCFE.east); % 情绪分支输出点
      \coordinate (cfe-domain-east) at ([yshift=-2.5mm]sharedCFE.east); % 域分支输出点
    }
  },
  % 情绪分类器组件
  pics/emotion classifier/.style={
    code={ % 定义情绪分类器
      \node[dann head] (emoHead) at (0, 0) {Emotion\\Classifier}; % 线性分类层
      \node[dann output] (emoOut) at (2.25, 0) {Predicted\\Label}; % 情感预测
      \draw[dann edge] (emoHead) -- (emoOut); % 分类器到输出
      \coordinate (emo-west) at (emoHead.west); % 左侧连接点
    }
  },
  % GRL组件
  pics/grl/.style={
    code={ % 定义梯度反转层
      \node[dann grl] (grlNode) at (0, 0) {GRL}; % 梯度反转层
      \coordinate (grl-west) at (grlNode.west); % 左侧连接点
      \coordinate (grl-east) at (grlNode.east); % 右侧连接点
    }
  },
  % 域判别器组件
  pics/domain discriminator/.style={
    code={ % 定义域判别器
      \node[dann domain] (domainHead) at (0, 0) {Domain\\Discriminator}; % 域判别器
      \node[dann output] (domainOut) at (2.45, 0) {Source / Target\\Confusion}; % 域混淆输出
      \draw[dann confuse edge] (domainHead) -- (domainOut); % 域判别到混淆
      \coordinate (domain-west) at (domainHead.west); % 左侧连接点
    }
  }
}

% 通用源域/目标域特征图例
\providecommand{\domainFeatureLegend}[2]{ % 锚点节点/下方距离
  \matrix (legendContent) [
    below=#2 of #1,
    column sep=2mm,
    inner sep=0pt,
    nodes={inner sep=0pt, anchor=center},
    ampersand replacement=\&
  ] {
    \node[minimum size=6mm] (legendSourceIcon) {}; \&
    \node[anchor=west, font=\scriptsize, text=gray!95] (legendSourceText) {Source feature}; \&
    \node[minimum width=5mm] {}; \&
    \node[minimum size=6mm] (legendTargetIcon) {}; \&
    \node[anchor=west, font=\scriptsize, text=gray!95] (legendTargetText) {Target feature}; \\
  }; % 图例内容矩阵
  \begin{scope}[shift={(legendSourceIcon.center)}, scale=.30, transform shape, name prefix=legendSource-]
    \pic {source feature cloud}; % 源域特征图例
  \end{scope}
  \begin{scope}[shift={(legendTargetIcon.center)}, scale=.30, transform shape, name prefix=legendTarget-]
    \pic {target feature cloud}; % 目标域特征图例
  \end{scope}
  \begin{scope}[on background layer]
    \node[
      rounded corners=4pt,
      draw=gray!85,
      dashed,
      inner xsep=5mm,
      inner ysep=2mm,
      fit=(legendContent)
    ] (legendBox) {}; % 图例外框
  \end{scope}
}
